% Auto-generated by build/ns_branch_successor.py. Do not edit by hand.

\section*{NS Branch Successor Pilot}

\addcontentsline{toc}{section}{NS Branch Successor Pilot}



\textbf{Primary governance function: stale-posture correction.} The NS governance repo has been computing \texttt{frontier-blocked} (from the stale inference rule \texttt{front:NS.main == O} $\Rightarrow$ \texttt{frontier-blocked}). The canonical NS branch book contains \texttt{prop:NS-status [C]}, \texttt{thm:NS71 [C]}, and \texttt{cor:NS-both-give-endpoint [C]} establishing \textbf{domain-bounded conditional CERT-CLOSE}. The NS preamble (``frontier-blocked'') is stale relative to the later sections of the same canonical file. This pilot corrects the posture and records the correction as a formal governance finding.



\subsection*{Stale-posture correction record}

Stale declared posture: \texttt{frontier-blocked}\\

Corrected computed posture: \textbf{\texttt{domain-bounded-cert-close}}\\

Basis: \texttt{thm:NS.7.1 [C]} $+$ \texttt{thm:NS.6.1 [C]} $+$ \texttt{front:NS.main == O}\\

Source: \texttt{prop:NS-status [C]} in canonical NS Branch Book



\subsection*{Domain-bounded closure hypotheses}

Phase-A, $K_0=7$, PASS, LS-2. Removing any hypothesis reopens the domain-bounded closure claim.



\subsection*{Active frontier}

\texttt{front:NS.main}: Stage-3 global unconditional closure (Navier-Stokes Clay Millennium Prize Problem). This is an \emph{external} frontier outside the current NFC toolkit. No admissible discharge mode exists within the current architecture.



\subsection*{SB2 and Ren.5 status}

\texttt{ob:NS.SB2}: conditionally discharged [C] via SCT.6b (\texttt{thm:NS.SB2-discharged}).\\

\texttt{thm:NS.Ren5}: conditionally discharged [C] via sub-lemma chain.



\begin{longtable}{p{0.30\linewidth}p{0.14\linewidth}p{0.10\linewidth}p{0.12\linewidth}p{0.20\linewidth}}
\caption{NS successor-pilot claim catalog}\label{tab:ns-successor-claims}\\
\textbf{ID} & \textbf{Kind} & \textbf{D.} & \textbf{C.} & \textbf{Scope}\\
\hline\endfirsthead
\textbf{ID} & \textbf{Kind} & \textbf{D.} & \textbf{C.} & \textbf{Scope}\\
\hline\endhead
thm:NS.stage0 & theorem & U & U & unconditional\\
thm:NS.stage1 & theorem & U & U & unconditional\\
thm:NS.stage2a & theorem & U & U & unconditional\\
thm:NS.stage2b & theorem & C & C & conditional-PASS\\
thm:NS.leray & theorem & U & U & unconditional\\
thm:NS.window-uniformity & theorem & C & C & conditional-Phase-A\\
thm:NS.6.1 & theorem & C & C & branch-local\\
thm:NS.Ren5 & theorem & C & C & conditional-Phase-A\\
thm:NS.SB2-discharged & theorem & C & C & conditional-SCT6b\\
ob:NS.SB2 & obligation & C & C & branch-local\\
ob:NS.bridge-stage2 & obligation & O & O & branch-local\\
thm:NS.7.1 & theorem & C & C & domain-bounded\\
cor:NS.domain-bounded-cert-close & corollary & C & C & domain-bounded\\
prop:NS.status & proposition & C & C & domain-bounded\\
front:NS.main & frontier & O & O & branch-local\\
status:NS & status\_note & domain-bounded-cert-close & domain-bounded-cert-close & branch-wide\\
\end{longtable}



\subsection*{Generated claim shells}

\begin{nfctheorem}{NS Stage-0 Thermodynamic Floor Density}
\textbf{ID.} \texttt{thm:NS.stage0}\quad\textbf{Decl.} U\quad\textbf{Comp.} U\\
\textbf{Scope.} regime=NS-source-descended; domain=NS-branch; certificate\_scope=unconditional; promotion\_ceiling=U

The thermodynamic floor density f\_inf > 0 exists unconditionally: under the NS active architecture declarations (N-add, T-add, BV, non-degeneracy), the Van Hove exhaustion window family has a positive limiting floor density. This is the first unconditional structural result of the NS branch and holds without any Phase-A or K\_0 hypothesis.
\end{nfctheorem}



\begin{nfctheorem}{NS Stage-1 Boundary Flux Stabilization}
\textbf{ID.} \texttt{thm:NS.stage1}\quad\textbf{Decl.} U\quad\textbf{Comp.} U\\
\textbf{Scope.} regime=NS-source-descended; domain=NS-branch; certificate\_scope=unconditional; promotion\_ceiling=U

The boundary flux phi\_inf >= 0 stabilizes unconditionally: the source-descended boundary flux of the NS observable family achieves a non-negative limit along the Van Hove exhaustion. This follows from the floor density (Stage-0) and the declared monotonicity of the transport regime.

\medskip
\noindent\textit{Dependencies.} thm:NS.stage0
\end{nfctheorem}



\begin{nfctheorem}{NS Stage-2a Discrete Balance Identity}
\textbf{ID.} \texttt{thm:NS.stage2a}\quad\textbf{Decl.} U\quad\textbf{Comp.} U\\
\textbf{Scope.} regime=NS-source-descended; domain=NS-branch; certificate\_scope=unconditional; promotion\_ceiling=U

The discrete balance identity holds unconditionally for the source-descended NS obstruction: for every test element f in the observable family, sum\_k <D\_k, f> = <D\_inf, f>. The identity is the canonical decomposition of the obstruction quantity into its transport and defect components. No continuum or PASS hypothesis is needed.

\medskip
\noindent\textit{Dependencies.} thm:NS.stage1
\end{nfctheorem}



\begin{nfctheorem}{NS Stage-2b Nonlinear Term Identification}
\textbf{ID.} \texttt{thm:NS.stage2b}\quad\textbf{Decl.} C\quad\textbf{Comp.} C\\
\textbf{Scope.} regime=NS-PASS-screened; domain=NS-branch; certificate\_scope=conditional-PASS; promotion\_ceiling=C; hyp=PASS-screening

The nonlinear term (u dot grad)u + grad p is identified from the discrete balance identity in the continuum limit, conditional on PASS screening: the transverse component of the nonlinear term vanishes under PASS, yielding the canonical Navier-Stokes nonlinear form. This identification requires the PASS condition from Book II and is therefore conditional [C].

\medskip
\noindent\textit{Hypotheses.} PASS-screening

\medskip
\noindent\textit{Dependencies.} thm:NS.stage2a
\end{nfctheorem}



\begin{nfctheorem}{NS Leray Weak Existence}
\textbf{ID.} \texttt{thm:NS.leray}\quad\textbf{Decl.} U\quad\textbf{Comp.} U\\
\textbf{Scope.} regime=NS-source-descended; domain=NS-branch; certificate\_scope=unconditional; promotion\_ceiling=U

For every initial datum u\_0 in L\textasciicircum{}2(R\textasciicircum{}3) with div u\_0 = 0, there exists a global Leray weak solution u in L\textasciicircum{}inf\_t L\textasciicircum{}2\_x intersect L\textasciicircum{}2\_t H\textasciicircum{}1\_x satisfying the energy inequality. The proof proceeds via Galerkin approximation on the NS observable family, using the discrete balance identity (Stage-2a), Aubin-Lions compactness, and PASS screening for incompressibility. This result is unconditional: no Phase-A or K\_0 hypothesis is needed. It is the canonical base for all subsequent NS closure arguments.

\medskip
\noindent\textit{Dependencies.} thm:NS.stage2a
\end{nfctheorem}



\begin{nfctheorem}{NS Window-Uniformity and Eventual Periodicity (Q2a Resolved)}
\textbf{ID.} \texttt{thm:NS.window-uniformity}\quad\textbf{Decl.} C\quad\textbf{Comp.} C\\
\textbf{Scope.} regime=NS-Phase-A-K0-7; domain=NS-branch; certificate\_scope=conditional-Phase-A; promotion\_ceiling=C; hyp=Phase-A,K0-7,kappa-bound

Under Phase-A, K\_0=7, and the kappa-bound (kappa <= 1 - 8/343), the NS active architecture window configurations eventually become periodic: there exists N\_0 such that for all n >= N\_0, the window configuration W\_n is determined by a finite state machine with periodic transitions. The uniform bounds c\_1 >= 1/7, c\_2 >= 1/49, eta <= 6/7, kappa <= 1 - 8/343 hold. This is the Q2a resolution: the NS active architecture declaration implies deterministic finite-state transition of window configurations, giving periodicity (not merely recurrence). Status: [C] conditional on Phase-A/K\_0=7.

\medskip
\noindent\textit{Hypotheses.} Phase-A, K0-7, kappa-bound

\medskip
\noindent\textit{Dependencies.} thm:NS.stage2a, thm:NS.6.1
\end{nfctheorem}



\begin{nfctheorem}{Active Missing Bridge Theorem}
\textbf{ID.} \texttt{thm:NS.6.1}\quad\textbf{Decl.} C\quad\textbf{Comp.} C\\
\textbf{Scope.} regime=declared-NS-target; domain=NS-branch; certificate\_scope=branch-local; promotion\_ceiling=C

Branch-local conditional bridge theorem placeholder for NS.

\medskip
\noindent\textit{Dependencies.} thm:VI.8.1, front:NS.main

\medskip
\noindent\textit{Discharge edges.} ob:NS.bridge-stage2 (partial)

\medskip
\noindent\textit{Proof stub.} Proof placeholder.
\end{nfctheorem}



\begin{nfctheorem}{Late Alignment-to-Dissipation Lemma (Ren.5)}
\textbf{ID.} \texttt{thm:NS.Ren5}\quad\textbf{Decl.} C\quad\textbf{Comp.} C\\
\textbf{Scope.} regime=NS-Phase-A-K0-7; domain=NS-branch; certificate\_scope=conditional-Phase-A; promotion\_ceiling=C; hyp=Phase-A,K0-7,PASS

The Late Alignment-to-Dissipation Lemma (Ren.5): under Phase-A, K\_0=7, and PASS, any infinite late counterexample to the alignment-to-dissipation property is eliminated by three-branch case analysis (sub-lemma chain NS.Ren.5a through NS.Ren.5e: Adaptation Dominance, Renewal-vs-Adaptation Trichotomy, Averaging Attenuation, Active Admissibility). All three branches being eliminated, Ren.5 holds inside the active architecture. The sub-lemma chain is conditional [C]. This discharges the proof-program frontier of the NS branch.

\medskip
\noindent\textit{Hypotheses.} Phase-A, K0-7, PASS

\medskip
\noindent\textit{Dependencies.} thm:NS.window-uniformity, thm:NS.6.1
\end{nfctheorem}



\begin{nfctheorem}{SB2 Conditionally Discharged via SCT.6b}
\textbf{ID.} \texttt{thm:NS.SB2-discharged}\quad\textbf{Decl.} C\quad\textbf{Comp.} C\\
\textbf{Scope.} regime=NS-Phase-A-K0-7; domain=NS-branch; certificate\_scope=conditional-SCT6b; promotion\_ceiling=C; hyp=Phase-A,K0-7,PASS,SCT.6b-firing-conditions

SB2 (Positive Reserve on the Safe Tail) is conditionally discharged via the SCT.6b coercive-transport mechanism: once a tail window satisfies the three SCT.6b firing conditions (coercive structure, transport compatibility, subcriticality check), SB2 is discharged from the assembled proof record via the shared coercive-transport doctrine (SCT.1-SCT.6b). The full discharge requires Ren.5 (thm:NS.Ren5) as the upstream input. Discharge mode: full, at the conditional scope of Phase-A/K\_0=7/PASS/SCT.6b.

\medskip
\noindent\textit{Hypotheses.} Phase-A, K0-7, PASS, SCT.6b-firing-conditions

\medskip
\noindent\textit{Dependencies.} thm:NS.Ren5

\medskip
\noindent\textit{Discharge edges.} ob:NS.SB2 (full)
\end{nfctheorem}



\begin{nfcobligation}{Positive Reserve on the Safe Tail}
\textbf{ID.} \texttt{ob:NS.SB2}\quad\textbf{Decl.} C\quad\textbf{Comp.} C\\
\textbf{Scope.} regime=declared-NS-target; domain=NS-branch; certificate\_scope=branch-local; promotion\_ceiling=C

Positive Reserve on the Safe Tail (SB2): CONDITIONALLY DISCHARGED via SCT.6b (thm:NS.SB2-discharged). The discharge is conditional on Phase-A/K\_0=7/PASS and the SCT.6b firing conditions. The obligation is closed at this scope. Unconditional force for SB2 would require O-IDcont-style force upgrade (external to current architecture).
\end{nfcobligation}



\begin{nfcobligation}{NS Bridge Stack Stage 2 Obligation}
\textbf{ID.} \texttt{ob:NS.bridge-stage2}\quad\textbf{Decl.} O\quad\textbf{Comp.} O\\
\textbf{Scope.} regime=declared-NS-target; domain=NS-branch; certificate\_scope=branch-local

Intermediate NS bridge-stage obligation used to test explicit discharge edges.

\medskip
\noindent\textit{Dependencies.} thm:VII.4.3
\end{nfcobligation}



\begin{nfctheorem}{Continuation-Criterion-to-Endpoint Theorem (NS.7.1)}
\textbf{ID.} \texttt{thm:NS.7.1}\quad\textbf{Decl.} C\quad\textbf{Comp.} C\\
\textbf{Scope.} regime=NS-Phase-A-K0-7-PASS-LS2; domain=NS-branch; certificate\_scope=domain-bounded; promotion\_ceiling=C; hyp=Phase-A,K0-7,PASS,LS-2,NS.6.1-discharged,Book-VI-H1x-interface

NS.7.1 (Continuation-Criterion-to-Endpoint Theorem): if the active NS persistence/extendibility continuation criterion holds in the promoted target scope (established by NS.6.1), then the NS endpoint target — the Leray weak solution u in L\textasciicircum{}inf\_t H\textasciicircum{}1\_x on the declared interval — follows. The proof combines: (1) Gronwall control from the multiplicative contraction D\_\{n+1\} <= kappa\textasciicircum{}\{n-n\_0\} D\_\{n\_0\} -> 0 (giving summable obstruction); (2) Book VI H\textasciicircum{}1\_x-norm legitimacy (licensing the smooth-notation statement). Domain-bounded CERT-CLOSE is achieved under Phase-A/K\_0=7/PASS/LS-2. NS.7.1 is structurally separate from NS.6.1 by canonical discipline: continuation and regularity are independently auditable.

\medskip
\noindent\textit{Hypotheses.} Phase-A, K0-7, PASS, LS-2, NS.6.1-discharged, Book-VI-H1x-interface

\medskip
\noindent\textit{Dependencies.} thm:NS.6.1, thm:NS.window-uniformity, thm:NS.SB2-discharged
\end{nfctheorem}



\begin{nfccorollary}{NS Domain-Bounded CERT-CLOSE}
\textbf{ID.} \texttt{cor:NS.domain-bounded-cert-close}\quad\textbf{Decl.} C\quad\textbf{Comp.} C\\
\textbf{Scope.} regime=NS-Phase-A-K0-7-PASS-LS2; domain=NS-branch; certificate\_scope=domain-bounded; promotion\_ceiling=C; hyp=Phase-A,K0-7,PASS,LS-2

The NS endpoint target (Leray weak solution u in L\textasciicircum{}inf\_t H\textasciicircum{}1\_x) follows from the declared source-descended obstruction-decay law, conditionally. The NS branch achieves domain-bounded CERT-CLOSE on the declared interval under Phase-A/K\_0=7/PASS/LS-2. Proof: Corollary of NS.6.1 (conditional) composed with NS.7.1 (conditional); both are [C]; their composition is [C]. Unconditional global regularity (Stage-3 Clay Prize) remains the external frontier.

\medskip
\noindent\textit{Hypotheses.} Phase-A, K0-7, PASS, LS-2

\medskip
\noindent\textit{Dependencies.} thm:NS.6.1, thm:NS.7.1
\end{nfccorollary}



\begin{nfcproposition}{NS Branch Status (Domain-Bounded Conditional CERT-CLOSE)}
\textbf{ID.} \texttt{prop:NS.status}\quad\textbf{Decl.} C\quad\textbf{Comp.} C\\
\textbf{Scope.} regime=NS-Phase-A-K0-7-PASS-LS2; domain=NS-branch; certificate\_scope=domain-bounded; promotion\_ceiling=C; hyp=Phase-A,K0-7,PASS,LS-2

NS Branch Status (Updated): The NS branch is a lawful descendant of the canonical spine and has reached domain-bounded conditional CERT-CLOSE under Phase-A/K\_0=7/PASS/LS-2 hypotheses. NS.6.1 is conditionally discharged (window-uniformity + K\_0=7 + Phase-A + PASS + kappa-bound). Ren.5 is conditionally discharged (sub-lemma chain NS.Ren.5a-5e). SB2 is conditionally discharged (via SCT.6b). NS.7.1 is conditionally discharged (Gronwall + Book VI H\textasciicircum{}1\_x interface). Domain-bounded CERT-CLOSE is achieved. The NS Frontier Theorem (thm:NS-frontier [U]) describes the formal conditional frontier; under the stated hypotheses the branch reaches closure. Unconditional global regularity (NS Stage-3 Clay Prize) remains the external frontier. Note: the NS preamble ("frontier-blocked") predates this discharge chain and is stale — the corrected posture is domain-bounded conditional CERT-CLOSE.

\medskip
\noindent\textit{Hypotheses.} Phase-A, K0-7, PASS, LS-2

\medskip
\noindent\textit{Dependencies.} cor:NS.domain-bounded-cert-close, thm:NS.SB2-discharged, thm:NS.Ren5
\end{nfcproposition}



\begin{nfcfrontier}{NS Main Frontier}
\textbf{ID.} \texttt{front:NS.main}\quad\textbf{Decl.} O\quad\textbf{Comp.} O\\
\textbf{Scope.} regime=declared-NS-target; domain=NS-branch; certificate\_scope=branch-local; promotion\_ceiling=O

Stage-3 global unconditional closure: prove unconditional global regularity for all smooth initial data for the 3D Navier-Stokes equations (the Clay Millennium Prize Problem). The domain-bounded conditional closure has been achieved under Phase-A/K\_0=7/PASS/LS-2 (thm:NS.7.1). This frontier records the residual gap between domain-bounded conditional CERT-CLOSE and unconditional global CERT-CLOSE. It is an external frontier: its resolution requires tools outside the current NFC active architecture. No admissible discharge mode exists within the current canonical toolkit.
\end{nfcfrontier}



\begin{nfcstatusnote}{NS Branch Posture Record}
\textbf{ID.} \texttt{status:NS}\quad\textbf{Decl.} domain-bounded-cert-close\quad\textbf{Comp.} domain-bounded-cert-close\\
\textbf{Scope.} regime=ns-branch-governance; domain=NS-branch; certificate\_scope=branch-wide; promotion\_ceiling=frontier-blocked

NS branch posture: domain-bounded-cert-close (conditionally). Correction note: the NS preamble says "frontier-blocked" but this is stale relative to the later sections of the canonical NS branch book, which establish domain-bounded conditional CERT-CLOSE via prop:NS.status [C]. The governance pilot records this stale-posture correction as a formal governance finding.

Active frontier: front:NS.main (Stage-3 global unconditional closure = Clay Prize). This is an external frontier outside the current NFC toolkit.

The domain-bounded closure is conditional on: Phase-A, K\_0=7, PASS, LS-2. Removing any of these hypotheses reopens the domain-bounded closure claim.

\medskip
\noindent\textit{Dependencies.} prop:NS.status, thm:V.UBLT, thm:VII.6.4
\end{nfcstatusnote}


